\documentclass[preamble=development-preamble.tex]{markdown}
title: The unary relational interpretation of arithmetic R
date: October 20, 2026

:::Abstract
The library `Interpretation_of_R.relational` gives a fixed interpretation
of the relational presentation of Tarski--Mostowski--Robinson arithmetic
in the original rank-based pure type system. The public theorem
`Theorem4.theorem4` combines realization of every axiom instance,
preservation of an independently defined classical deduction system,
proofs and internal refutations for all ground atoms, and a concrete
failure of reflection for arbitrary sentences. The construction uses a
unary Church carrier, negative equality tests, successor-closed
observations, and uniformly certified finite selectors. This exposition
follows the checked source declarations, including the fixed-rank
invariant behind the quantified axiom schemes.
:::

# The theorem and its source language

The source theory, written $\widetilde R$, has zero, successor, equality,
order, and ternary addition and multiplication relations $A,M$. Its
terms are variables, zero, and iterated successor; addition and
multiplication are predicates rather than source term constructors.
For external natural numbers $m,n$, its axiom schemes are

$$
\begin{aligned}
 A(\bar m,\bar n,x)&\leftrightarrow \overline{m+n}=x,\\
 M(\bar m,\bar n,x)&\leftrightarrow \overline{mn}=x,\\
 \bar m&\ne\bar n\quad(m\ne n),\\
 x\leq\bar n&\leftrightarrow\bigvee_{i=0}^{n}x=\bar i,\\
 x\leq\bar n&\lor\bar n\leq x.
\end{aligned}
$$

Free $x$ is universally quantified. Both directions of the bounded-order
biconditional are part of this relational presentation.
`DirectRRelationalSyntax.v` defines the scoped source syntax and these
five axiom constructors. Primitive logical syntax consists of falsity,
implication, and universal quantification; negation, conjunction,
biconditional, finite disjunction, and existential quantification are
defined encodings. The inductive judgment `Relational.derives` has
hypothesis and theory rules, implication and universal rules, classical
double-negation elimination, equality reflexivity, and equality
replacement. Its definition contains no target typing or inhabitation
premise.

Write $A^*$ for the fixed translation of a source sentence. The target
judgment below is exactly `PTSrk.Typing.typ` at `Ranks.RankSorts`, using
the original annotated `Syntax.Term nat n`.

:::Theorem[The relational interpretation]
The theorem `Theorem4.theorem4` packages the following four conclusions.

1. Every closed source axiom $A$ has a closed original PTS inhabitant
   $p:A^*$.
2. If `Relational.derives [] A`, there is a closed original PTS inhabitant
   $p:A^*$.
3. Every true ground equality, comparison, addition atom, or multiplication
   atom has a proof of its translation; every false ground atom has an
   internal refutation of type $A^*\to\bot$.
4. Some source sentence has an inhabited translation but no derivation in
   the independently defined source calculus.
:::

This is the theorem labeled `main` in the harmonized manuscript. Its
number is Theorem 4 there; it is distinct from the functional theorem and
from the functional `Kernel.Theorem5Target` interface. The conjunction
above has no assumed certificate package, all-axiom premise, or deduction
preservation premise. The subsidiary files establish these facts before
the final theorem assembles them.

# The fixed carrier, domain, and logical translation

The PTS has sorts $U_i:U_0$, and products have rank
$\max(a+1,b)$ when their domain and codomain have ranks $a,b$.
There is no cumulativity or eta conversion. Define the unary Church carrier

$$
 D:=\Pi X:U_1.X\to(X\to X)\to X:U_3,
 \qquad z:=\lambda X\,x\,s.x,
 \qquad S d:=\lambda X\,x\,s.s(d\,X\,x\,s).
$$

Displayed lambdas abbreviate fully annotated original terms. Let
$\nu_n=S^n z$. The module `Unary` also defines two auxiliary operations

$$
\begin{aligned}
 \operatorname{Add}(d,e)&:=\lambda X\,x\,s.d\,X\,(e\,X\,x\,s)\,s,\\
 \operatorname{Mul}(d,e)&:=\lambda X\,x\,s.
 d\,X\,x\,(\lambda y:X.e\,X\,y\,s).
\end{aligned}
$$

Both have type $D\to D\to D$. They appear inside translated relations;
they do not extend the source term syntax. The theorems
`UnaryEvaluation.addition_numerals` and `multiplication_numerals` prove
their ground beta equations. No equation identifying an arbitrary Church
presentation with its eta expansion is needed.

The admitted domain is a constructor-closure certificate:

$$
 \operatorname{Good}(d):=
 \Pi Q:D\to U_{10}.Q(z)\to
 (\Pi u:D.Q(u)\to Q(Su))\to Q(d):U_{12}.
$$

`Unary.good_zero_typed` and `good_successor_typed` certify its zero and
successor constructors. `admitted_numeral_typed` produces the certificates
for all external numeral indices. This definition supplies the precise
bounded-rank induction needed later; it is not an induction rule for
every raw $D$-inhabitant and does not assume the desired arithmetic axioms.

The fixed answer type is $\bot:=\Pi Z:U_0.Z$, and
$\top:=U_0\to U_0$ has its identity proof. Negation is $P\to\bot$,
conjunction is $(P\to Q\to\bot)\to\bot$, and a finite disjunction is
the flat continuation type

$$
 \bigvee_c(P_0,\ldots,P_k):=
 (P_0\to\bot)\to\cdots\to(P_k\to\bot)\to\bot.
$$

A proof type is stable when its double negation implies it.
The translation preserves implication and interprets universal
quantification by

$$
 (\forall x.B)^*:=\Pi d:D.\operatorname{Good}(d)\to B^*(d).
$$

Existential quantification follows from its source encoding and hence has
a negative, guarded translation. `RelationalTranslation.formula_stable`
proves stability for every translated formula. Classical elimination is
therefore realized by PTS terms; no elimination from $\bot$ into arbitrary
positive types is postulated.

# Equality and order through observations

## Negative equality tests

For each $P:D\to U_1$, compare the negative tests $P(d)\to\bot$ and
$P(e)\to\bot$. With the encoded stable biconditional, define

$$
 E(d,e):=\Pi P:D\to U_1.
       \bigl(\neg_c P(d)\leftrightarrow_c\neg_c P(e)\bigr):U_5.
$$

Equality reflexivity, symmetry, and transitivity follow pointwise by
identity, reversal, and composition of implications. Every typed
$F:D\to D$ respects $E$: a test on $F(d)$ is pulled back to the test
$u\mapsto P(F(u))$ on $d$. `Unary.equality_congruence_typed` implements
that argument. It yields congruence for successor and both arguments of
the auxiliary arithmetic operations.

Consequently the arithmetic predicates are

$$
 A^*(d,e,f):=E(\operatorname{Add}(d,e),f),
 \qquad M^*(d,e,f):=E(\operatorname{Mul}(d,e),f).
$$

Their equality compatibility follows from operation congruence and the
equivalence laws. Equality is thus a definable logical relation; it is
not asserted to coincide with target beta conversion.

## Successor-closed observation tests

A unary observation consists of the dependent telescope

$$
 \mathcal U=(X:U_1)(x_0:X)(s:X\to X)(\delta:X\to U_4)
 (d_0:\delta(x_0))
 (d_s:\Pi u:X.\delta(u)\to\delta(su)).
$$

For $d:D$, its observation is
$\operatorname{ev}_{\mathcal U}(d)=d\,X\,x_0\,s$.
Define the positive predicate

$$
 G(d):=\Pi\mathcal U.\delta(\operatorname{ev}_{\mathcal U}(d)):U_6.
$$

For a test $P:X\to U_1$, its upward certificate says

$$
 \operatorname{Up}(P):=
 \Pi u:X.\delta(u)\to\neg_c P(u)\to\neg_c P(su).
$$

Order transports a negative test through every certified observation:

$$
 O(d,e):=\Pi\mathcal U.\Pi P:X\to U_1.
 \operatorname{Up}(P)\to
 \neg_c P(\operatorname{ev}_{\mathcal U}(d))\to
 \neg_c P(\operatorname{ev}_{\mathcal U}(e)):U_6.
$$

The observation quantifies only the displayed initial, successor, and
upward certificates. Reflexivity and transitivity of $O$ are derived by
identity and composition; they are not extra observation premises.
Specializing equality to the pulled-back test
$u\mapsto P(\operatorname{ev}_{\mathcal U}(u))$ proves $E(d,e)\to O(d,e)$.
Symmetry of $E$ and transitivity then give congruence of order.

The beta fold equations establish zero and successor closure of $G$.
To derive `Unary.good_observed_typed`, use the rank-ten motive
$u\mapsto J_9\to G(u)$ in the `Good` certificate and evaluate at the
canonical inhabitant of $J_9$. This explicit rank raising replaces an
unavailable cumulativity step. Finally $G(d)$ supplies the guard required
to derive $O(d,Sd)$ from the upward certificate. Repeating this step gives
`Unary.ground_order_realized` for every external $i\leq j$.

# Finite separation: turning a wrong comparison into Bottom

:::Theorem[No backward numeral step]
For every external $n$, there is an original inhabitant of
$O(\nu_{n+1},\nu_n)\to\bot$. The declaration is `no_back_realized`
in the `Separation` module.
:::

:::Proof
Set $c=n+1$ and $q=c+1$. Use the selector carrier
$F_q=U_1\to\cdots\to U_1\to U_1$ with selectors $e_0,\ldots,e_c$,
zero $e_0$, and the capped successor table
$s(e_i)=_\beta e_{\min(i+1,c)}$. The shared table-action and iteration
theorems prove that the numeral $\nu_k$ evaluates to
$e_{\min(k,c)}$ in this observation.

Its positive domain predicate is

$$
 \delta_q(u):=\Pi Q:F_q\to U_3.
 Q(e_0)\to\cdots\to Q(e_c)\to Q(u):U_4.
$$

Selector projection supplies the initial certificate. For successor
closure, specialize an input certificate to $v\mapsto Q(sv)$ after
opening the desired output predicate $Q$. Each resulting branch reduces
to a selector branch already supplied by the continuation.

Take the separating test

$$
 P(u):=u\,\underbrace{\top\cdots\top}_{c\text{ entries}}\,\bot.
$$

The upward law is proved by membership elimination at the rank-three
motive $v\mapsto\neg_c P(v)\to\neg_c P(sv)$. At the cap it is identity.
Below the cap its premise is $\top\to\bot$, which yields $\bot$ when
applied to the known truth inhabitant; abstraction supplies the negative
conclusion. These branches construct the required upward certificate.

Now specialize a purported $O(\nu_c,\nu_n)$ proof to this complete,
closed observation. It gives
$(\bot\to\bot)\to(\top\to\bot)$. Applying the two available identity
proofs yields $\bot$.
:::

`Separation.ground_order_refutation` extends this result to every reversed
ground comparison, using a forward comparison from $n+1$ to the larger
index. `Separation.R3_realized` refutes equality of any two distinct
numerals by $E\Rightarrow O$ and, when necessary, symmetry of $E$.
These finite observations have internal positive certificates; no
parametricity principle claiming that all raw selectors are canonical is
used. They are separating observations, not finite models of every
source axiom.

# Bounded cases and the quantified schemes

The key all-index theorem classifies an admitted presentation with a
global order alternative:

$$
 \operatorname{Case}_n(d):=
 \bigvee_c\bigl(E(d,\nu_0),\ldots,E(d,\nu_n),O(\nu_{n+1},d)\bigr):U_8.
$$

The last alternative is an actual global lower bound. Merely observing
overflow in one finite algebra would not provide it.

:::Theorem[Bounded cases]
For every external cutoff $n$, `Unary.bounded_case_realized` constructs a
closed original proof of
$\Pi d:D.\operatorname{Good}(d)\to\operatorname{Case}_n(d)$.
:::

:::Proof
Use the positive rank-ten invariant

$$
 \operatorname{Inv}_n(d):=
 \Pi Z:U_8.(G(d)\to\operatorname{Case}_n(d)\to Z)\to Z.
$$

At zero, supply its positive observation certificate and the reflexive
equality branch. At successor, open the parent's continuation and
eliminate the child invariant at that same $Z$. This supplies both
$G(d)$ and the child case proof. Construct $G(Sd)$ from $G(d)$
independently. To construct the new negative case proof, open all branch
continuations, so that the remaining answer type is $\bot$.

An equality branch with index $i<n$ moves to $i+1$ by successor
congruence. At $i=n$, successor congruence and equality-to-order give
$O(\nu_{n+1},Sd)$. The already-large branch composes its lower bound
with $O(d,Sd)$, whose positive guard is the retained $G(d)$.
This proves successor closure of the invariant without extracting
positive evidence from a negative disjunction.

Apply the original `Good` certificate to this exact rank-ten invariant.
Project its case component by taking $Z=\operatorname{Case}_n(d)$ and
the continuation $\lambda g\,c.c$. The source proves formation of the
entire quantified case type at rank thirteen as well.
:::

The bounded-order scheme now follows in both directions.
Under $O(d,\nu_n)$, a finite branch gives the desired finite equality
disjunction. The large branch composes to $O(\nu_{n+1},\nu_n)$,
contradicting separation. Conversely, an $i\leq n$ equality branch gives
$O(d,\nu_n)$ through equality-to-order and the ground comparison.
`UnaryBounded.R4_realized` assembles the biconditional.

For comparability, a finite branch gives $O(d,\nu_n)$, while a large
branch composed with $O(\nu_n,\nu_{n+1})$ gives $O(\nu_n,d)$.
`UnaryAxioms.R5_realized` supplies the resulting disjunction. The first
two graph schemes follow from the auxiliary arithmetic beta equations
for arbitrary admitted third arguments. Numeral inequality is the
separation result. `UnaryArithmetic.all_axioms_realized` and
`RelationalAxioms.source_axiom_realized` connect these constructions to
the literal source axiom syntax.

# Preserving deductions and characterizing ground truth

The translation uses an environment from source variables to original
PTS terms. `term_typed` and `term_admitted` prove that source terms remain
in the admitted carrier when their environment entries do. Source
renaming and substitution commute with translation; in particular,
`instantiate_translation` identifies the translated universal instance
with target substitution.

`RelationalEquality.formula_transport` extends atomic congruence to all
formulas and `equality_replacement` realizes the source replacement rule.
The important restriction is that source variables range over admitted
arguments. The proof does not assume or require that `Good` is invariant
under equality between arbitrary raw presentations.

`RelationalDeduction.preservation` is proved by induction on the
independent source derivation. Its open-context premises say that
environment entries are typed and admitted and that translated source
hypotheses are inhabited. Universal introduction extends the target
context with both a carrier variable and its domain proof. Universal
elimination uses term admission and the substitution identity. Classical
elimination uses formula stability, and equality replacement uses the
transport theorem. The closed specialization is
`RelationalDeduction.sentence_preservation`.

For ground truth, a closed source term has a natural-number value and
translates to the corresponding unary numeral up to beta conversion.
Equality and order are then handled by reflexivity, forward comparisons,
and separation. Addition and multiplication atoms also use the two
arithmetic beta equations. `DirectRRelationalGround.v` proves each true
and false case, and `Theorem4.ground_correctness` packages all four atom
forms. These are proofs and internal refutations for every ground atom,
not a finite regression sample and not reflection of arbitrary sentences.

# Why the interpretation is not faithful

Consider $\rho:=\forall x.(x\leq x)$. Its translation is inhabited by
`Unary.admitted_order_reflexivity`: after opening the domain and
observation binders, the proof returns its negative-test argument
unchanged. Order is reflexive even on raw presentations, so the domain
proof is not needed by this term.

For source non-derivability, the shared module `SourceCountermodel` uses
$\mathbb N\sqcup\{\omega\}$. Successor fixes $\omega$; the arithmetic
graphs have the usual values on natural inputs and value zero when an
input is $\omega$. Order is usual on naturals, each natural is below
$\omega$, and $\omega$ is below no point. Thus every standard numeral
has exactly its expected finite downset, and every point is comparable
with a standard numeral, but $\omega\not\leq\omega$.

`RelationalCountermodel` defines evaluation and satisfaction for the
independent source syntax, proves each axiom, and proves deduction
soundness. Satisfaction is stable by structural induction: atomic
relations in this model are decidable, and implication and universal
quantification preserve stability. This justifies the classical source
rule constructively, without a native excluded-middle axiom.
`order_reflexivity_underivable` therefore gives

$$
 \widetilde R\nvdash\rho,
 \qquad\text{while}\qquad
 \vdash\operatorname{orderRefl}:\rho^*.
$$

`Theorem4.nonfaithfulness` packages this witness. The countermodel is a
model of the source arithmetic, not a set-theoretic model of the PTS.
Its exceptional point also explains why nonfaithfulness does not conflict
with the complete ground proof/refutation construction.

# Reading map, verification, and limits

`DirectRUnary.v` contains the fixed terms, rank and typing certificates,
basic equality/order laws, and bounded-case invariant. The evaluation,
separation, bounded, arithmetic, and unary-axiom modules build the five
schemes. The `DirectRRelational*` modules then supply independent syntax,
structural translation, formula transport, deduction preservation,
ground results, and countermodel soundness. `DirectRTheorem4.v` is the
public endpoint. Shared logic, selectors, and certificate construction
are documented by the parent development's exposition.

The central ranks do not depend on an axiom's numeral indices or a
finite observation's size:

| Construction | Rank |
| --- | --- |
| Unary carrier $D$ | $3$ |
| Finite domain certificate $\delta_q$ | $4$ |
| Equality and arithmetic graph atoms | $5$ |
| Positive observations $G$ and order $O$ | $6$ |
| Bounded cases | $8$ |
| Positive invariant | $10$ |
| Admitted domain predicate | $12$ |

These are checked sufficient ranks; no optimality theorem is claimed.
The construction provides neither induction for all raw carrier terms nor
reflection of arbitrary source derivability. Internal refutations have
type $A\to\bot$; external uninhabitability of $\bot$ is a separate
metatheoretic assertion, not a conjunct of Theorem 4.

The library name remains `Interpretation_of_R.relational`. To inspect its
public endpoint in Rocq:

```coq
From Interpretation_of_R.relational Require Import DirectRTheorem4.
Check Theorem4.theorem4.
Print Assumptions Theorem4.theorem4.
```

With the pinned environment described in `docs/DirectR-CI.md`, build from
the repository root:

```sh
dune build @theories/Interpretation_of_R/relational/audit
```

The alias includes the unary and relational audit modules and depends
only on this interpretation, the shared library, and the original
kernel. It uses OCaml 4.14.2, Dune 3.23.1, Rocq 9.2.0, Stdlib 9.1.0, and
Equations 1.3.2+9.2. Use a fresh absolute Dune build directory when fresh
assumption output is required.

The retained 14 October layout audit prints the exact public theorem
above followed by “Closed under the global context”, and records
successful fresh compilation and recursive kernel checking. Its
`COMMANDS.md` and `types.log` are under
`docs/direct-r-handoff/verification/theories-layout-2026-10-14`.
The loaded-context inventory's unchanged upstream functional
extensionality is distinct from the theorem's closed transitive
assumption report. More recent verification is recorded in the live
handoff. `bash tools/check-local-ci.sh` runs the complete staged local
audit, including pristine upstream; compiling this exposition alone is
not a fresh Rocq audit.
